% ==============================================================================
% PART III: DISCOURSE LOGIC & ARGUMENTATION ENGINE
% File: chapters/part3_discourse.tex
% ==============================================================================
\part{Discourse Logic \& Argumentation Engine / 篇章逻辑与论证引擎}\label{part:discourse}

\chapter{The Information Flow Paradigm / 篇章信息流与主位推进模型}\label{ch:infoflow}
优秀的学术段落不是由离散句子的简单拼盘，而是一条精密咬合的认知齿轮链条。读者在理解学术技术论述时，依赖每一句话句首对前文概念的自然锚定。本章系统阐述主位推进（Theme-Rheme Progression）的底层规律，确立以已知信息引导新信息的根本法则。

\section{The Given-to-New Fundamental Law (L001) / 旧信息在前、新信息在后原则}
句子包含两个核心信息层级：主位（Theme，通常位于句首，代表已知信息 Given）与述位（Rheme，位于句中和句末，代表新信息 New）。连贯段落的核心秘诀在于：后一句的主位必须承接前一句的述位或已知概念，形成链式信息流。

\section{Thematic Progression Architectures (L002 \& L003) / 两大主流主位推进架构}
\begin{enumerate}[label=\textbf{\color{AcadNavy}\arabic*.}, leftmargin=2em]
  \item \textbf{恒定主位推进（Constant Theme, L002）}：多句话保持相同主位实体（如同一种算法架构），分别在述位抛出其在速度、精度、内存与鲁棒性等多个维度的评价。
  \item \textbf{线性链式推进（Linear Rheme-to-Theme, L003）}：第一句句尾的新概念 $Y$，立即作为第二句句首的已知主位锚点，进而引出新概念 $Z$，形成 $X \to Y \Rightarrow Y \to Z$ 的因果推导闭环。
\end{enumerate}

\section{Anaphoric Anchoring \& Cohesive Scaffolding (L004 \& L005) / 壳名词指代与语义钩子}
严禁使用裸露的代词（如单字的 \textit{This shows that...}）。应当使用具有概括总结功能的“壳名词（Shell Nouns）”，如 \textit{This discrepancy, This trade-off, Such an anomaly}，既指代明确，又赋予前文定性评价。

\chapter{Toulmin Micro-Argumentation Architecture / 图尔敏微论证架构与科学思辨}\label{ch:toulmin}
科学论证不是情绪宣泄，而是基于证据的严密辩护。现代学术话语体系广泛采用图尔敏论证模型（Toulmin Model of Argumentation）来组织段落内的微论证，确保每个学术主张均有实证基准与逻辑保证的支撑。

\section{The Six-Constituent Toulmin Formalism (L006) / 图尔敏六要素论证模型}
\begin{itemize}[leftmargin=1.5em]
  \item \textbf{主张 (Claim)}: 本段的核心学术论点或结论。
  \item \textbf{数据/事实 (Data)}: 实验指标、统计观测数据或实证记录。
  \item \textbf{保证 (Warrant)}: 连接数据与主张的科学定律、理论桥梁或算法先验。
  \item \textbf{支撑 (Backing)}: 证明保证本身成立的底层数学基础或前人基石文献。
  \item \textbf{反驳/例外 (Rebuttal)}: 坦诚承认主张在某些极端边界条件下可能受限。
  \item \textbf{限制词 (Qualifier)}: 审慎限定主张成立的概率范围（typically, substantially, in finite samples）。
\end{itemize}

\section{Mechanistic Causality vs. Spurious Correlation (L007) / 机制因果与伪相关辨析}
在实验分析段落中，严禁仅凭统计相关性直接断言因果关系。必须交代物理媒介或内部特征流的传导过程，彻底杜绝逻辑假飞跃。

\section{Concession, Pivot, and Structural Refutation (L008) / 让步、转折与结构化反驳}
在学术商榷与审稿回复中，标准的绅士论辩采用“承认合理 $\to$ 聚焦边界 $\to$ 扭转权重”三步舞姿：首先诚恳承认竞品或审稿人顾虑的合理性，继而指出其隐含假设的脆弱性，最后证明本文方案在关键工况下的决定性收益。

\chapter{Zero AI-ese Discourse Surgery \& 30 Academic Logic Rules (L001--L030) / 去 AI 化篇章外科重构与 30 条逻辑法则}\label{ch:logic30}
随着大语言模型的普及，审稿人对充斥着 \textit{delve into, paramount, testaments, tapestry, paves the way} 等浮夸套话的 AI-ese 风格表现出极度反感。本章设立去 AI 化篇章重构手术准则，并完整解析 30 条学术逻辑与论证核心法则（L001--L030）。

\section{The Zero AI-ese Surgical Protocol (L029) / 去 AI 化外科重构准则}
坚决割除大模型生成的“四大浮夸病灶”：
\begin{enumerate}[label=\textbf{\color{TargetHighlight}Symptom \arabic*:}, leftmargin=2em]
  \item \textbf{空洞修辞帽}: 诸如 \textit{In the ever-evolving landscape of...}，一律删除，直奔领域实体。
  \item \textbf{过度吹捧}: 诸如 \textit{groundbreaking, unprecedented, game-changer}，替换为具体的量化突破指标。
  \item \textbf{散装总结}: 诸如 \textit{In conclusion, it is worth noting that...}，替换为紧凑的因果或对比收束。
  \item \textbf{动词贫乏}: 诸如大量滥用 \textit{plays a crucial role in}，替换为具体动作及物动词（underpins, mitigates, regulates）。
\end{enumerate}

\section{The Complete 30 Academic Logic Rules Encyclopedia (L001--L030) / 30条学术逻辑法则全集}
\Needspace{6\baselineskip}
\subsection{L001: Given-to-New Information Flow / 旧信息在前、新信息在后原则（Given-to-New Principle）}\label{sec:logic:L001}
\begin{ConceptBox}{Core Logic Principle \& Structural Formula}
\textbf{Category}: \texttt{Information Flow (信息流推进)}\par\vspace{3pt}
\textbf{Core Principle}: 英语段落连贯的第一动力不是连接词，而是每一句的句首（主位 Theme）承接上一句句末已经出现过的概念（旧信息 Given），再在句末（述位 Rheme）抛出本句真正要推进的核心增量（新信息 New）。\par\vspace{4pt}
\textbf{Structural Formula}:\par\noindent\texttt{Sentence A: [Subject X] $\to$ [Verb] $\to$ [New Concept Y]. | Sentence B: [This/Such Y or Pronoun] $\to$ [Verb] $\to$ [New Concept Z].}
\end{ConceptBox}
\begin{ExampleBox}{Defective Draft vs. Idiomatic Reconstruction}
\textbf{\color{TargetHighlight}[Defective Draft]}:\par\textit{Autoregressive transformers suffer from severe error accumulation over long horizons. A sudden distribution shift in non-stationary sensor streams further amplifies this error accumulation.}\par\vspace{4pt}
\textbf{\color{AcadNavy}[Idiomatic Reconstruction]}:\par\textit{Autoregressive transformers suffer from severe error accumulation over long horizons. This cumulative drift accelerates sharply when non-stationary sensor streams undergo sudden distribution shifts.}\par\vspace{4pt}
\textbf{Surgical Analysis}: 修改前第二句以长新名词 A sudden distribution shift... 开头，旧信息 this error accumulation 被甩在句末；修改后第二句用 This cumulative drift 紧扣上一句句尾的 error accumulation，再引出新信息 distribution shifts，段落瞬间丝滑。
\end{ExampleBox}
\begin{TipBox}{Academic Trap \& Reviewer Perspective}
\textbf{Common Pitfall}: 每一句都突兀地以一个全新概念作长主语，而把指代上一句的词塞在句尾，导致读者每读一句都要重新建立思维锚点。
\end{TipBox}
\vspace{6pt}
\subsection{L002: Constant Theme Progression / 平行主位推进模式（Constant Theme Progression）}\label{sec:logic:L002}
\begin{ConceptBox}{Core Logic Principle \& Structural Formula}
\textbf{Category}: \texttt{Information Flow (信息流推进)}\par\vspace{3pt}
\textbf{Core Principle}: 当一段话旨在从多个互补维度刻画同一个核心模型、算法或社会现象时，连续 2--3 句保持相同或同指代的主语（如 Our framework / The temporal encoder / It），让谓语动词承担信息递进。\par\vspace{4pt}
\textbf{Structural Formula}:\par\noindent\texttt{S1: [Theme T] + [Action 1 (Architecture)]. | S2: [T / It] + [Action 2 (Optimization)]. | S3: [By combining both, T] + [Action 3 (Empirical Outcome)].}
\end{ConceptBox}
\begin{ExampleBox}{Defective Draft vs. Idiomatic Reconstruction}
\textbf{\color{TargetHighlight}[Defective Draft]}:\par\textit{The encoder extracts multi-scale patches first. Then frequency spectrum filtering is performed by a Fourier block. Finally, high-resolution predictions are generated by the decoder.}\par\vspace{4pt}
\textbf{\color{AcadNavy}[Idiomatic Reconstruction]}:\par\textit{Our architecture processes non-stationary sequences in three stages. It first partitions the raw series into multi-scale temporal patches, then filters high-frequency noise in the Fourier domain, and finally decodes these refined representations into multi-horizon forecasts.}\par\vspace{4pt}
\textbf{Surgical Analysis}: 修改前主语在 encoder、filtering、predictions 之间无序跳转且混用被动语态；修改后以 Our architecture / It 为统一主位，用三个并列强动词 partitions $\to$ filters $\to$ decodes 清晰串联数据流。
\end{ExampleBox}
\begin{TipBox}{Academic Trap \& Reviewer Perspective}
\textbf{Common Pitfall}: 为了所谓“句式多变”强行每句换主语、频繁在主动与被动之间横跳，导致读者搞不清到底是谁在执行模块动作。
\end{TipBox}
\vspace{6pt}
\subsection{L003: Linear Rheme-to-Theme Chain / 延续链式推进模式（Linear Chain Progression）}\label{sec:logic:L003}
\begin{ConceptBox}{Core Logic Principle \& Structural Formula}
\textbf{Category}: \texttt{Information Flow (信息流推进)}\par\vspace{3pt}
\textbf{Core Principle}: 阐述因果机制、算法推导或仿真到真实迁移（Sim-to-Real）链条时，严格采用 A $\to$ B，B $\to$ C，C $\to$ D 的咬合结构，上一句的宾语转化为下一句的主语。\par\vspace{4pt}
\textbf{Structural Formula}:\par\noindent\texttt{S1: A induces B. $\to$ S2: This B distorts C. $\to$ S3: Once C degrades, D collapses.}
\end{ConceptBox}
\begin{ExampleBox}{Defective Draft vs. Idiomatic Reconstruction}
\textbf{\color{TargetHighlight}[Defective Draft]}:\par\textit{Unmodeled friction in simulated joints causes real-world grasping failure. Domain randomization often fails because state transitions diverge.}\par\vspace{4pt}
\textbf{\color{AcadNavy}[Idiomatic Reconstruction]}:\par\textit{Unmodeled joint friction in simulators creates a systematic discrepancy between synthetic and physical state transitions. This transition mismatch misleads visuomotor policies during contact-rich phases, ultimately causing real-world grasping failures.}\par\vspace{4pt}
\textbf{Surgical Analysis}: 通过 synthetic and physical state transitions $\to$ This transition mismatch $\to$ contact-rich phases $\to$ grasping failures 形成严密的四环机械传导链。
\end{ExampleBox}
\begin{TipBox}{Academic Trap \& Reviewer Perspective}
\textbf{Common Pitfall}: 中间跳过关键传导环节（如直接从 A 跳到 D），或者链条顺序颠倒。
\end{TipBox}
\vspace{6pt}
\subsection{L004: Shell Noun / Summary Noun Anaphora / 概念浓缩指代词（This/Such + Summary Noun）}\label{sec:logic:L004}
\begin{ConceptBox}{Core Logic Principle \& Structural Formula}
\textbf{Category}: \texttt{Cohesion \& Reference (指代与衔接)}\par\vspace{3pt}
\textbf{Core Principle}: 严禁单独使用裸露的 This / That / It 指代上一整句话！必须在指示代词后紧跟一个精准概括上一句本质的抽象范畴名词（如 This discrepancy / Such degradation / These computational constraints / This trade-off）。\par\vspace{4pt}
\textbf{Structural Formula}:\par\noindent\texttt{Sentence 1: [Full empirical observation or limitation]. $\to$ Sentence 2: This/Such [Summary Noun: bottleneck / discrepancy / vulnerability / shift] + [Active Verb: motivates / necessitates / limits / explains]...}
\end{ConceptBox}
\begin{ExampleBox}{Defective Draft vs. Idiomatic Reconstruction}
\textbf{\color{TargetHighlight}[Defective Draft]}:\par\textit{Diffusion models require hundreds of iterative denoising steps at inference time. This prevents them from running on onboard robot controllers.}\par\vspace{4pt}
\textbf{\color{AcadNavy}[Idiomatic Reconstruction]}:\par\textit{Diffusion policies require hundreds of iterative denoising steps at inference time. This computational latency prevents them from running at closed-loop control frequencies on onboard robot hardware.}\par\vspace{4pt}
\textbf{Surgical Analysis}: 用 This computational latency 将上一句“数百步去噪迭代”精准定性为“计算延迟”，既消除了指代歧义，又完成了因果定性。
\end{ExampleBox}
\begin{TipBox}{Academic Trap \& Reviewer Perspective}
\textbf{Common Pitfall}: 句首单写 'This makes the model slow' 或 'Which causes error'，审稿人和雅思考官无法判断 This 究竟指上一句的主语、宾语还是整件事。
\end{TipBox}
\vspace{6pt}
\subsection{L005: Semantic Hooks over Mechanical Connectors / 去机械连接词化：用语义钩子替代 First/Second/Moreover}\label{sec:logic:L005}
\begin{ConceptBox}{Core Logic Principle \& Structural Formula}
\textbf{Category}: \texttt{Cohesion \& Reference (指代与衔接)}\par\vspace{3pt}
\textbf{Core Principle}: 母语学者与高分作文极少每句句首都贴 Firstly, Secondly, Moreover, What is more, In addition 等膏药。应使用语义承接状语（Beyond accuracy gains, / Alongside temporal drift, / At the hardware level,）实现自然递进。\par\vspace{4pt}
\textbf{Structural Formula}:\par\noindent\texttt{Instead of 'Moreover, X does Y', use: 'Beyond [Summary of Prior Point], [Subject X] also [Verb Y].'}
\end{ConceptBox}
\begin{ExampleBox}{Defective Draft vs. Idiomatic Reconstruction}
\textbf{\color{TargetHighlight}[Defective Draft]}:\par\textit{First, our model reduces MSE by 15\%. Second, our model uses 40\% less memory. Moreover, our model generalizes to unseen sensors.}\par\vspace{4pt}
\textbf{\color{AcadNavy}[Idiomatic Reconstruction]}:\par\textit{Our model reduces forecasting MSE by 15.4\% across six benchmarks. Beyond accuracy gains, its sparse attention mechanism cuts peak GPU memory by 40\%, enabling zero-shot deployment on unseen sensor arrays.}\par\vspace{4pt}
\textbf{Surgical Analysis}: 用 Beyond accuracy gains 承上启下，用 enabling... 自然引出第三层泛化收益，彻底摆脱机械序数词。
\end{ExampleBox}
\begin{TipBox}{Academic Trap \& Reviewer Perspective}
\textbf{Common Pitfall}: 一段 5 句话连续用 First, Second, Besides, Moreover，而句子内部毫无词汇呼应，形成典型的中式“八股排比腔”或 AI 列表腔。
\end{TipBox}
\vspace{6pt}
\subsection{L006: Claim-Mechanism-Evidence-Warrant-Qualification / 图尔敏五步微论证模型（Toulmin Micro-Argumentation）}\label{sec:logic:L006}
\begin{ConceptBox}{Core Logic Principle \& Structural Formula}
\textbf{Category}: \texttt{Argumentation (论证结构)}\par\vspace{3pt}
\textbf{Core Principle}: 任何一个独立成段的学术或议论主张，都必须具备完整的五要素：(1)  主张（Claim）、(2)  机理解释（Mechanism/Reason）、(3)  实证数据或具象案例（Evidence）、(4)  逻辑桥梁（Warrant：为什么该证据能证明该主张）、(5)  边界限定（Qualification/Hedging）。\par\vspace{4pt}
\textbf{Structural Formula}:\par\noindent\texttt{[S1: Claim] $\to$ [S2: Underlying Mechanism] $\to$ [S3: Concrete Metric/Instance] $\to$ [S4: Warrant/Implication] $\to$ [S5: Boundary/Scope Qualification].}
\end{ConceptBox}
\begin{ExampleBox}{Defective Draft vs. Idiomatic Reconstruction}
\textbf{\color{TargetHighlight}[Defective Draft]}:\par\textit{Patch-level tokenization is much better than point-wise tokenization. It captures local information very well and improves time series forecasting significantly in all cases.}\par\vspace{4pt}
\textbf{\color{AcadNavy}[Idiomatic Reconstruction]}:\par\textit{Patch-level tokenization outperforms point-wise representations in long-horizon time series forecasting [Claim]. Whereas single time steps lack semantic context and amplify high-frequency noise, contiguous sub-series patches preserve local slope and periodicity [Mechanism]. On the ETTm2 benchmark, grouping 16 consecutive steps reduces 720-step MSE by 18.7\% compared with point-wise attention [Evidence]. This gain confirms that local structural aggregation mitigates attention dispersion over long lookback windows [Warrant], although excessively large patches may blur abrupt transient anomalies [Qualification].}\par\vspace{4pt}
\textbf{Surgical Analysis}: 五句话各司其职：立论 $\to$ 对比机理 $\to$ ETTm2 实证数据 $\to$ 理论印证 $\to$ 诚实指出过大 Patch 的边界局限，无懈可击。
\end{ExampleBox}
\begin{TipBox}{Academic Trap \& Reviewer Perspective}
\textbf{Common Pitfall}: 只有 Claim + 换句话重复一遍 Claim + 空洞夸张形容词，中间既无传导机制也无具象证据。
\end{TipBox}
\vspace{6pt}
\subsection{L007: Causal Mechanism vs. Spurious Correlation / 因果机制链 vs. 表面相关性切分（Causality vs. Correlation）}\label{sec:logic:L007}
\begin{ConceptBox}{Core Logic Principle \& Structural Formula}
\textbf{Category}: \texttt{Argumentation (论证结构)}\par\vspace{3pt}
\textbf{Core Principle}: 在论文实验分析与雅思/托福议论文中，观察到 A 与 B 同时变化时，绝不能直接断言 A causes B；必须区分“直接因果（drives / induces）”、“中介传导（mediates / channels）”与“混杂相关（coincides with / correlates with）”。\par\vspace{4pt}
\textbf{Structural Formula}:\par\noindent\texttt{Observation: A correlates with B. $\to$ Confounder check: Rather than A directly causing B, confounding factor C drives both $\to$ Controlled verification.}
\end{ConceptBox}
\begin{ExampleBox}{Defective Draft vs. Idiomatic Reconstruction}
\textbf{\color{TargetHighlight}[Defective Draft]}:\par\textit{Increasing model parameters makes the robot grasp better in the real world.}\par\vspace{4pt}
\textbf{\color{AcadNavy}[Idiomatic Reconstruction]}:\par\textit{While scaling backbone parameters correlates with higher success rates in simulation, real-world grasping gains stem primarily from diverse tactile pre-training rather than parameter count alone.}\par\vspace{4pt}
\textbf{Surgical Analysis}: 精准用 correlates with 界定表面伴随关系，再用 stem primarily from... rather than... 点明真实因果驱动力。
\end{ExampleBox}
\begin{TipBox}{Academic Trap \& Reviewer Perspective}
\textbf{Common Pitfall}: 将伴随发生的数据现象直接写成绝对因果，或用 make / let 等口语词粗糙表达因果。
\end{TipBox}
\vspace{6pt}
\subsection{L008: Concession, Pivot, and Structural Refutation / 先让步、再釜底抽薪式反驳（Concession-Pivot-Refutation）}\label{sec:logic:L008}
\begin{ConceptBox}{Core Logic Principle \& Structural Formula}
\textbf{Category}: \texttt{Argumentation (论证结构)}\par\vspace{3pt}
\textbf{Core Principle}: 无论是论文引言批评前人方法、Rebuttal 回复审稿人，还是六级/雅思大作文回应反方观点，最高级的说服力都来自“先承认对方在特定受限条件下确实有效（Concession），再指出该条件在真实/复杂场景下不成立（Pivot \& Refutation）”。\par\vspace{4pt}
\textbf{Structural Formula}:\par\noindent\texttt{Admittedly / While [Prior Method / Opposing View] achieves [Specific Merit] under [Strict Assumption A], this assumption breaks down when [Real-World Condition B], resulting in [Concrete Failure].}
\end{ConceptBox}
\begin{ExampleBox}{Defective Draft vs. Idiomatic Reconstruction}
\textbf{\color{TargetHighlight}[Defective Draft]}:\par\textit{MLP-based forecasters are too simple and cannot model complex dependencies, so our Transformer model is necessary.}\par\vspace{4pt}
\textbf{\color{AcadNavy}[Idiomatic Reconstruction]}:\par\textit{While lightweight MLP forecasters achieve competitive accuracy on stationary sequences with stable periodicity, their fixed temporal weights cannot adapt dynamically when industrial sensor streams undergo abrupt regime shifts.}\par\vspace{4pt}
\textbf{Surgical Analysis}: 先肯定 MLP 在平稳周期序列上的竞争力，再精准指出其“固定时间权重无法适应突发工况漂移”的结构性软肋，有理有据。
\end{ExampleBox}
\begin{TipBox}{Academic Trap \& Reviewer Perspective}
\textbf{Common Pitfall}: 一上来全盘否定前人工作（Previous works are totally wrong/useless），显得极不客观且容易激怒审稿人。
\end{TipBox}
\vspace{6pt}
\subsection{L009: Taxonomic Synthesis vs. Chronological Laundry List / 分类矩阵综述法：拒绝 A did X, B did Y 流水账（Taxonomy over Laundry List）}\label{sec:logic:L009}
\begin{ConceptBox}{Core Logic Principle \& Structural Formula}
\textbf{Category}: \texttt{Paragraph Architecture (段落架构)}\par\vspace{3pt}
\textbf{Core Principle}: 撰写 Related Work 或文献背景段时，严禁以人名作主语逐句罗列（Smith (2022) proposed X. Lee (2023) proposed Y. Wang (2024) used Z.）。必须以“方法流派/技术范式”为主语进行聚类，把文献引用退居句尾括号或作同位语例证。\par\vspace{4pt}
\textbf{Structural Formula}:\par\noindent\texttt{Existing approaches to [Problem] fall into two primary paradigms: [Paradigm A] and [Paradigm B]. [Paradigm A] methods (e.g., [1, 2]) rely on..., whereas [Paradigm B] architectures [3, 4] leverage... However, neither paradigm resolves [Core Bottleneck].}
\end{ConceptBox}
\begin{ExampleBox}{Defective Draft vs. Idiomatic Reconstruction}
\textbf{\color{TargetHighlight}[Defective Draft]}:\par\textit{Wu et al. proposed Autoformer in 2021. Zhou et al. proposed FEDformer in 2022. Nie et al. proposed PatchTST in 2023. They all improved time series forecasting.}\par\vspace{4pt}
\textbf{\color{AcadNavy}[Idiomatic Reconstruction]}:\par\textit{Recent deep forecasters mitigate long-range attention complexity either by decomposing seasonal-trend components in the frequency domain (e.g., Autoformer, FEDformer) or by aggregating local sub-series into channel-independent patches (e.g., PatchTST). Yet both paradigms treat cross-variable correlations as static over time, overlooking transient inter-sensor coupling during system faults.}\par\vspace{4pt}
\textbf{Surgical Analysis}: 两句话完成流派二分法归纳 + 代表作锚定 + 共同盲区提炼，学术密度提升数倍。
\end{ExampleBox}
\begin{TipBox}{Academic Trap \& Reviewer Perspective}
\textbf{Common Pitfall}: 连续 6 句话全部是 [Author] et al. proposed/introduced/developed...，只有罗列没有批判性归类。
\end{TipBox}
\vspace{6pt}
\subsection{L010: The Metric-Baseline-Condition Precision Triad / 定量实证三元组：指标 + 基线 + 边界条件（Metric-Baseline-Condition Triad）}\label{sec:logic:L010}
\begin{ConceptBox}{Core Logic Principle \& Structural Formula}
\textbf{Category}: \texttt{Precision \& Grounding (精确度与实证落地)}\par\vspace{3pt}
\textbf{Core Principle}: 任何宣称“更好、更快、更鲁棒、影响深远”的句子，必须同时交代三个维度：(1)  具体评估指标与幅度（Metric \& Magnitude）、(2)  对比参照系（Baseline）、(3)  评测工况或数据集条件（Condition）。\par\vspace{4pt}
\textbf{Structural Formula}:\par\noindent\texttt{[Method/Policy] improves/reduces [Metric] by [Exact \%/Value] compared with [Strongest Baseline] under [Specific Dataset/Horizon/Noise Regime].}
\end{ConceptBox}
\begin{ExampleBox}{Defective Draft vs. Idiomatic Reconstruction}
\textbf{\color{TargetHighlight}[Defective Draft]}:\par\textit{Our generative policy achieves much better success rates and runs much faster than previous methods.}\par\vspace{4pt}
\textbf{\color{AcadNavy}[Idiomatic Reconstruction]}:\par\textit{Our consistency-distilled policy increases bimanual peg-insertion success rates from 64\% to 89\% while cutting inference latency from 180 ms to 12 ms relative to a 50-step DDPM baseline.}\par\vspace{4pt}
\textbf{Surgical Analysis}: 用 64\%$\to$89\%（指标）、180ms$\to$12ms（延迟）、relative to a 50-step DDPM baseline（基线）和 bimanual peg-insertion（任务条件）彻底消灭空洞形容词。
\end{ExampleBox}
\begin{TipBox}{Academic Trap \& Reviewer Perspective}
\textbf{Common Pitfall}: 只写 significantly improves performance 或 achieves state-of-the-art results，属于典型的“不精确 + 空洞夸张”。
\end{TipBox}
\vspace{6pt}
\subsection{L011: Subject Compression and End-Weight Balance / 段落内主语瘦身与尾重平衡（Subject Compression \& End-Weight）}\label{sec:logic:L011}
\begin{ConceptBox}{Core Logic Principle \& Structural Formula}
\textbf{Category}: \texttt{Sentence \& Paragraph Rhythm (节奏与尾重)}\par\vspace{3pt}
\textbf{Core Principle}: 在多句段落中，严禁出现超过 10--12 个单词的臃肿主语从句或连环 of 短语作主语（头重脚轻）。把已知的短核心名词放在主语位置（≤6 词），将复杂的条件、机制和结果顺延到谓语动词之后的宾语或状语从句中。\par\vspace{4pt}
\textbf{Structural Formula}:\par\noindent\texttt{Bad: [The phenomenon that A with B under C fails to D] is common. $\to$ Good: [Short Subject A] frequently fails to [Verb D] when [Condition B and C occur].}
\end{ConceptBox}
\begin{ExampleBox}{Defective Draft vs. Idiomatic Reconstruction}
\textbf{\color{TargetHighlight}[Defective Draft]}:\par\textit{The phenomenon that autoregressive language models hallucinate unverified citations when prompted with out-of-distribution academic queries is a serious obstacle.}\par\vspace{4pt}
\textbf{\color{AcadNavy}[Idiomatic Reconstruction]}:\par\textit{Autoregressive language models frequently hallucinate unverified citations when prompted with out-of-distribution academic queries, posing a serious obstacle to automated literature review.}\par\vspace{4pt}
\textbf{Surgical Analysis}: 主语从 17 个词压缩为 3 个词（Autoregressive language models），强动词 hallucinate 提前，句尾用分词结果状语自然收束。
\end{ExampleBox}
\begin{TipBox}{Academic Trap \& Reviewer Perspective}
\textbf{Common Pitfall}: 受中文“……的现象/问题/事实表明……”影响，写出 The fact that... is... 或 The problem of the degradation of... is... 的畸形句。
\end{TipBox}
\vspace{6pt}
\subsection{L012: One Paragraph, One Controlling Idea / 一段一核：单一控制中心法则（Single Controlling Idea Rule）}\label{sec:logic:L012}
\begin{ConceptBox}{Core Logic Principle \& Structural Formula}
\textbf{Category}: \texttt{Paragraph Architecture (段落架构)}\par\vspace{3pt}
\textbf{Core Principle}: 一个自然段只能回答一个核心问题（例如：这一段只讲“为什么频域滤波能抑制传感器高频噪声”，就绝不要在同一段末尾顺手塞进“此外我们还开源了代码并降低了显存”）。\par\vspace{4pt}
\textbf{Structural Formula}:\par\noindent\texttt{Topic Sentence (Controlling Idea X) $\to$ Mechanism of X $\to$ Evidence for X $\to$ Boundary/Synthesis of X. (Move Concept Y to Next Paragraph).}
\end{ConceptBox}
\begin{ExampleBox}{Defective Draft vs. Idiomatic Reconstruction}
\textbf{\color{TargetHighlight}[Defective Draft]}:\par\textit{Our spatial-temporal attention block captures cross-sensor correlations accurately. Also, training takes only two hours on a single GPU, and the public dataset we used has some missing timestamps that we interpolated linearly.}\par\vspace{4pt}
\textbf{\color{AcadNavy}[Idiomatic Reconstruction]}:\par\textit{Our spatial-temporal attention block captures dynamic cross-sensor correlations by weighting neighbor nodes according to instantaneous phase lag. This phase-aware aggregation prevents faulty sensors from corrupting adjacent channels during localized hardware spikes, reducing multi-sensor MAE by 11.8\%.}\par\vspace{4pt}
\textbf{Surgical Analysis}: 砍掉无关的训练耗时与数据清洗杂音，全段聚焦于“相移感知图注意力如何隔离故障传感器并提升多传感器预测精度”。
\end{ExampleBox}
\begin{TipBox}{Academic Trap \& Reviewer Perspective}
\textbf{Common Pitfall}: 一段话里既讲精度、又讲速度、还顺带抱怨数据集质量，导致段落像杂货铺一样失焦。
\end{TipBox}
\vspace{6pt}
\subsection{L013: Deductive Zoom-In vs. Inductive Zoom-Out / 演绎推进（Zoom-In）与归纳提炼（Zoom-Out）双模切换}\label{sec:logic:L013}
\begin{ConceptBox}{Core Logic Principle \& Structural Formula}
\textbf{Category}: \texttt{Paragraph Architecture (段落架构)}\par\vspace{3pt}
\textbf{Core Principle}: 写方法论与议论文主论点段宜用演绎法（Zoom-In：首句抛出概括性规律，随后逐句拆解机理与实例）；写实验分析与看图小作文局部段宜用归纳法（Zoom-Out：先报告具体观测对比，末句提炼背后的理论规律）。\par\vspace{4pt}
\textbf{Structural Formula}:\par\noindent\texttt{Zoom-In: [General Claim] $\to$ [Specific Mechanism] $\to$ [Concrete Data]. | Zoom-Out: [Data Observation A vs. B] $\to$ [Pattern Across Horizons] $\to$ [Theoretical Takeaway].}
\end{ConceptBox}
\begin{ExampleBox}{Defective Draft vs. Idiomatic Reconstruction}
\textbf{\color{TargetHighlight}[Defective Draft]}:\par\textit{At horizon 96, our MSE is 0.182. At horizon 336, our MSE is 0.245. At horizon 720, our MSE is 0.310, while Informer reaches 0.612.}\par\vspace{4pt}
\textbf{\color{AcadNavy}[Idiomatic Reconstruction]}:\par\textit{As the prediction horizon extends from 96 to 720 steps, our model's MSE rises modestly from 0.182 to 0.310, whereas Informer degrades sharply to 0.612. This widening performance gap demonstrates that decoupling trend projection from seasonal attention effectively suppresses long-range error compounding.}\par\vspace{4pt}
\textbf{Surgical Analysis}: 第一句整合多步长数据趋势对比，第二句用 This widening performance gap demonstrates that... 完成归纳升华，点明“趋势与季节解耦抑制长程误差累积”的本质原因。
\end{ExampleBox}
\begin{TipBox}{Academic Trap \& Reviewer Perspective}
\textbf{Common Pitfall}: 在实验分析段只罗列完数据就戛然而止，缺少最后一句 Zoom-Out 的机理升华（Why this pattern matters）。
\end{TipBox}
\vspace{6pt}
\subsection{L014: Intuition-First Equation Integration / 公式与散文无缝咬合：先给直觉动机，后推符号定义（Intuition-Before-Math）}\label{sec:logic:L014}
\begin{ConceptBox}{Core Logic Principle \& Structural Formula}
\textbf{Category}: \texttt{Methodology Writing (方法章节逻辑)}\par\vspace{3pt}
\textbf{Core Principle}: 在 Methodology 段落中引入数学公式时，严禁突兀甩出 'The equation is defined as follows:'。必须遵循三步曲：(1)  物理/算法直觉（为什么要算这个量） $\to$ (2)  数学定义（把公式作为句子的语法宾语或同位语自然嵌入） $\to$ (3)  关键符号的物理含义与边界。\par\vspace{4pt}
\textbf{Structural Formula}:\par\noindent\texttt{To [Intuitive Goal], we formulate [Concept] as: [Equation], where [Symbol A] denotes [Meaning] and [Symbol B] constrains [Boundary].}
\end{ConceptBox}
\begin{ExampleBox}{Defective Draft vs. Idiomatic Reconstruction}
\textbf{\color{TargetHighlight}[Defective Draft]}:\par\textit{We use a gating weight. The formula is w = sigmoid(Wx + b). In this formula, x is the feature and w is the weight.}\par\vspace{4pt}
\textbf{\color{AcadNavy}[Idiomatic Reconstruction]}:\par\textit{To suppress high-frequency sensor jitter before temporal aggregation, we compute a channel-wise attenuation gate $g_t = \sigma(W_g h_t + b_g)$, where $h_t \in \mathbb{R}^d$ denotes the instantaneous hidden state and $g_t \in (0, 1)^d$ dynamically scales down noisy channels.}\par\vspace{4pt}
\textbf{Surgical Analysis}: 一句话之内完成“目的状语（抑制高频抖动）+ 核心主谓宾（计算通道衰减门）+ 公式同位语 + where 定语从句解释符号维度与物理作用”，紧凑专业。
\end{ExampleBox}
\begin{TipBox}{Academic Trap \& Reviewer Perspective}
\textbf{Common Pitfall}: 公式前后全是碎片句（Here is the formula: ... In the formula, x is input, y is output.），破坏段落连贯。
\end{TipBox}
\vspace{6pt}
\subsection{L015: Pattern Interpretation over Table Recitation / 图表叙事铁律：解读规律而非朗读表格（Interpret Patterns, Never Read Tables Aloud）}\label{sec:logic:L015}
\begin{ConceptBox}{Core Logic Principle \& Structural Formula}
\textbf{Category}: \texttt{Experimental Writing (实验章节逻辑)}\par\vspace{3pt}
\textbf{Core Principle}: 读者自己会看表格数字。实验段落的任务不是把 Table 1 的每一行数字抄进正文，而是回答三个问题：(1)  跨数据集的整体胜率与极值在哪？(2)  哪类基线在什么条件下崩盘了？(3)  这种差异验证了我们在引言/方法中提出的什么假设？\par\vspace{4pt}
\textbf{Structural Formula}:\par\noindent\texttt{[Table Reference + Overall Dominance] $\to$ [Contrast between Baseline Families under Hard Regime] $\to$ [Mechanistic Explanation linking back to Method].}
\end{ConceptBox}
\begin{ExampleBox}{Defective Draft vs. Idiomatic Reconstruction}
\textbf{\color{TargetHighlight}[Defective Draft]}:\par\textit{Table 1 shows the results. On dataset A, our model gets 91\%. On dataset B, our model gets 88\%. Our model is the best on all datasets.}\par\vspace{4pt}
\textbf{\color{AcadNavy}[Idiomatic Reconstruction]}:\par\textit{Table 1 shows that our model achieves the lowest forecasting error on 26 out of 30 benchmark settings. While point-wise Transformers degrade rapidly on volatile electricity load datasets (e.g., ECL), our multi-scale spectral mixer reduces MSE by an average of 16.4\%, confirming that frequency-domain denoising stabilizes long-horizon representations.}\par\vspace{4pt}
\textbf{Surgical Analysis}: 既有全局概览（26/30 最优），又有困难场景（波动剧烈的电力负载 ECL）下的流派对比，最后回扣方法论核心假设。
\end{ExampleBox}
\begin{TipBox}{Academic Trap \& Reviewer Perspective}
\textbf{Common Pitfall}: 把表格每一行数字机械抄进正文，只报数不提炼规律与原因。
\end{TipBox}
\vspace{6pt}
\subsection{L016: Counterfactual Ablation Attribution / 消融实验归因逻辑：移除模块 $\to$ 暴露失效模式 $\to$ 反证模块必要性（Ablation Attribution）}\label{sec:logic:L016}
\begin{ConceptBox}{Core Logic Principle \& Structural Formula}
\textbf{Category}: \texttt{Experimental Writing (实验章节逻辑)}\par\vspace{3pt}
\textbf{Core Principle}: 消融实验段落（Ablation Study）的核心逻辑是“反事实归因”：不仅要报告拿掉模块 A 后指标下降了多少，更要指出拿掉 A 后模型具体在什么子任务/子场景上出现了什么退化，从而锁定模块 A 的不可替代性。\par\vspace{4pt}
\textbf{Structural Formula}:\par\noindent\texttt{Removing [Module A] (w/o A) causes [Metric] to drop by [X\%], with the sharpest degradation occurring under [Specific Hard Condition]. This drop indicates that without [Module A], the model fails to [Specific Function].}
\end{ConceptBox}
\begin{ExampleBox}{Defective Draft vs. Idiomatic Reconstruction}
\textbf{\color{TargetHighlight}[Defective Draft]}:\par\textit{When we remove the kinematics constraint loss, the success rate drops by 14\%. This shows that the kinematics constraint loss is very important for our robot.}\par\vspace{4pt}
\textbf{\color{AcadNavy}[Idiomatic Reconstruction]}:\par\textit{Removing the kinematic feasibility loss (w/o $\mathcal{L}_{\text{kin}}$) reduces real-world drawer-opening success from 86\% to 72\%. Inspection of failed trajectories reveals that without $\mathcal{L}_{\text{kin}}$, the diffusion policy frequently generates wrist orientations that exceed joint torque limits during the pull phase.}\par\vspace{4pt}
\textbf{Surgical Analysis}: 不仅给出 86\%$\to$72\% 的量化回落，还通过轨迹检查揭示了“缺少运动学约束导致手腕姿态超出关节扭矩极限”的具体物理死因。
\end{ExampleBox}
\begin{TipBox}{Academic Trap \& Reviewer Perspective}
\textbf{Common Pitfall}: 只写 Without module A, accuracy drops by 3\%, which proves module A is very important and effective（同义反复，未解释为什么降）。
\end{TipBox}
\vspace{6pt}
\subsection{L017: Constructive Boundary Framing in Limitations / 局限性（Limitations）的无损诚实写法：界定适用边界而非自毁贡献}\label{sec:logic:L017}
\begin{ConceptBox}{Core Logic Principle \& Structural Formula}
\textbf{Category}: \texttt{Discussion \& Hedging (边界与审慎表达)}\par\vspace{3pt}
\textbf{Core Principle}: 顶会强制要求写 Limitations。高水平写法是：明确指出当前方法在什么极端假设违背下（如极端稀疏采样、毫秒级超高速碰撞）性能会受限，并立即给出造成该边界的客观算力/物理原因及可行的下一步演进路径。\par\vspace{4pt}
\textbf{Structural Formula}:\par\noindent\texttt{While our framework performs robustly under [Standard Regime], its accuracy/efficiency is currently bounded when [Extreme Condition]. This limitation arises because [Objective Mechanism], suggesting a promising direction in [Concrete Future Fix].}
\end{ConceptBox}
\begin{ExampleBox}{Defective Draft vs. Idiomatic Reconstruction}
\textbf{\color{TargetHighlight}[Defective Draft]}:\par\textit{Our model has some limitations. It is too slow when there are many variables and does not work well on irregular time series.}\par\vspace{4pt}
\textbf{\color{AcadNavy}[Idiomatic Reconstruction]}:\par\textit{While our cross-variate graph module scales linearly with sequence length, its quadratic complexity with respect to channel count ($O(C^2)$) increases memory overhead when $C > 1,000$. Incorporating low-rank channel projection or sparse k-NN graph routing offers a direct path to scaling toward ultra-high-dimensional sensor arrays.}\par\vspace{4pt}
\textbf{Surgical Analysis}: 把“慢”精确化为通道维度的 $O(C^2)$ 复杂度边界（$C > 1,000$），并顺势给出低秩投影或稀疏 k-NN 路由的建设性解法。
\end{ExampleBox}
\begin{TipBox}{Academic Trap \& Reviewer Perspective}
\textbf{Common Pitfall}: 要么写成虚假自夸（Our only limitation is that our model is too general），要么写成致命自毁（Our model fails completely on real data）。
\end{TipBox}
\vspace{6pt}
\subsection{L018: Spoken Rhythm, Visual Anchoring, and Rhetorical Signposting / 口头学术演讲的听觉信息流：短分句、视觉锚点与设问引导（Oral Signposting）}\label{sec:logic:L018}
\begin{ConceptBox}{Core Logic Principle \& Structural Formula}
\textbf{Category}: \texttt{Presentation \& Speech (学术演讲与汇报)}\par\vspace{3pt}
\textbf{Core Principle}: 听众无法像读论文那样倒回去重读长定语从句。把书面段转化为演讲稿时，必须：(1)  砍掉多重嵌套从句，拆为 10--15 词的清晰短句；(2)  先给视觉坐标（Look at the red curve in the upper-right panel）再讲结论；(3)  用轻量设问（Why does this spike happen? Because...）引导注意力。\par\vspace{4pt}
\textbf{Structural Formula}:\par\noindent\texttt{[Visual Anchor: Turning to the right panel of Slide 4,] $\to$ [Short Rhetorical Question: why does the baseline diverge after step 50?] $\to$ [Crisp Causal Answer: The answer lies in compound orientation drift.]}
\end{ConceptBox}
\begin{ExampleBox}{Defective Draft vs. Idiomatic Reconstruction}
\textbf{\color{TargetHighlight}[Defective Draft]}:\par\textit{As illustrated in Figure 4, the trajectory generated by the baseline model without temporal consistency regularization exhibits substantial oscillation around the target waypoint.}\par\vspace{4pt}
\textbf{\color{AcadNavy}[Idiomatic Reconstruction]}:\par\textit{Take a look at the dashed blue trajectory in Figure 4. Notice how violently the robot arm oscillates near the target waypoint. Why does this happen? Without temporal consistency regularization, adjacent action chunks conflict with one another.}\par\vspace{4pt}
\textbf{Surgical Analysis}: 用 Take a look at... $\to$ Notice how... $\to$ Why does this happen? $\to$ Without X, Y conflicts... 营造极具现场掌控力的母语级 Tech Talk 节奏。
\end{ExampleBox}
\begin{TipBox}{Academic Trap \& Reviewer Perspective}
\textbf{Common Pitfall}: 拿着书面论文原封不动照念长达 40 个词的复合句，台下听众在前 10 秒就丢失了主谓关系。
\end{TipBox}
\vspace{6pt}
\subsection{L019: Evidence-First Polite Rebuttal Protocol / 顶会 Rebuttal 黄金四步：致谢认同 $\to$ 结论前置 $\to$ 新增证据/表格 $\to$ 修订位置锚定}\label{sec:logic:L019}
\begin{ConceptBox}{Core Logic Principle \& Structural Formula}
\textbf{Category}: \texttt{Rebuttal \& Peer Review (审稿回复逻辑)}\par\vspace{3pt}
\textbf{Core Principle}: 回复审稿人尖锐质疑时，绝不能情绪化辩解或绕圈子。第一句简短感谢并重申问题核心；第二句直接给出明确结论（Direct Answer）；第三至四句呈上新增实验数据或理论澄清（New Empirical Evidence）；末句指明已在修订稿哪一节哪一表更新。\par\vspace{4pt}
\textbf{Structural Formula}:\par\noindent\texttt{We thank Reviewer X for raising this insightful question regarding [Topic]. [Direct Answer: Our method maintains 84\% success even under 30\% occlusion]. Specifically, our new ablation in Table R1 shows that [Data]. We have incorporated this analysis into Section 4.3 (highlighted in blue).}
\end{ConceptBox}
\begin{ExampleBox}{Defective Draft vs. Idiomatic Reconstruction}
\textbf{\color{TargetHighlight}[Defective Draft]}:\par\textit{You misunderstood our paper. We already explained this in Section 3. Our model is actually very robust to noise.}\par\vspace{4pt}
\textbf{\color{AcadNavy}[Idiomatic Reconstruction]}:\par\textit{We thank the reviewer for highlighting the question of sensor noise robustness. To verify performance under severe corruption, we evaluated our policy under Gaussian observation noise ($\sigma \in [0.05, 0.20]$). As shown in Table R2 below, our model degrades by only 3.2\% at $\sigma=0.20$, compared with a 21.5\% drop for the Diffusion Policy baseline. We have added these results to Appendix D.}\par\vspace{4pt}
\textbf{Surgical Analysis}: 礼貌专业，直接用新增噪声区间实验数据（3.2\% vs. 21.5\% 跌幅对比）让审稿人心服口服。
\end{ExampleBox}
\begin{TipBox}{Academic Trap \& Reviewer Perspective}
\textbf{Common Pitfall}: 用大段文字教育审稿人“你没看懂我们的论文”，或者只口头承诺“我们以后会做实验”而不给真实数据。
\end{TipBox}
\vspace{6pt}
\subsection{L020: Counterfactual Consequence Chain / 反事实假设论证法（Counterfactual Reasoning in Exam Essays）}\label{sec:logic:L020}
\begin{ConceptBox}{Core Logic Principle \& Structural Formula}
\textbf{Category}: \texttt{Exam Essays (六级/雅思/托福论证)}\par\vspace{3pt}
\textbf{Core Principle}: 在六级、雅思 Task 2 和托福独立写作中，当正面说理写满两句感觉词穷时，立即引入一句“反事实推演（Absent this policy / Without such regulation / Were universities to abandon X）”：指出如果缺失该条件，会引发什么具体负面连锁反应，从而反向锁死中心论点。\par\vspace{4pt}
\textbf{Structural Formula}:\par\noindent\texttt{[Positive Mechanism S1-S2]. Conversely, without [Core Factor X], [System/Individuals] would inevitably face [Concrete Negative Outcome Y], as seen when [Real-World Scenario Z].}
\end{ConceptBox}
\begin{ExampleBox}{Defective Draft vs. Idiomatic Reconstruction}
\textbf{\color{TargetHighlight}[Defective Draft]}:\par\textit{Independent critical thinking is very important for college students in the AI era. Students should think independently so that they can have a bright future.}\par\vspace{4pt}
\textbf{\color{AcadNavy}[Idiomatic Reconstruction]}:\par\textit{Training students to verify algorithmic outputs against primary sources cultivates rigorous analytical judgment. Without this verification habit, undergraduates risk accepting plausible-sounding AI hallucinations as factual premises, eroding their capacity to diagnose complex engineering failures.}\par\vspace{4pt}
\textbf{Surgical Analysis}: 第一句正面阐述“核查一手文献培养分析判断力”，第二句用 Without this verification habit... 反推不核查带来的具体后果（把逼真的 AI 幻觉当成事实前提），论证立体深刻。
\end{ExampleBox}
\begin{TipBox}{Academic Trap \& Reviewer Perspective}
\textbf{Common Pitfall}: 正面车轱辘话来回重复三遍（X is important. Therefore we must do X because X is very beneficial to society.）。
\end{TipBox}
\vspace{6pt}
\subsection{L021: Concrete Single-Sentence Micro-Case Illustration / 具象微观场景举例法：用一句话精确案例终结空洞模板（Micro-Case Grounding）}\label{sec:logic:L021}
\begin{ConceptBox}{Core Logic Principle \& Structural Formula}
\textbf{Category}: \texttt{Exam Essays (六级/雅思/托福论证)}\par\vspace{3pt}
\textbf{Core Principle}: 六级/雅思/托福作文严禁编造虚假的调查数据（如 *A recent survey by Harvard shows 87.3\% of people...*），也严禁写空洞口号（*For example, technology helps doctors a lot*）。最高分做法是用一句 `For instance, when [Specific Actor] implements [Specific Tool/Policy] in [Specific Scenario], [Concrete Measurable Effect]` 构建具象行业微观切片。\par\vspace{4pt}
\textbf{Structural Formula}:\par\noindent\texttt{For instance, when [municipal grid operators / radiology departments / logistics warehouses] deploy [specific technology/policy], they can [concrete operational improvement], thereby [link back to paragraph claim].}
\end{ConceptBox}
\begin{ExampleBox}{Defective Draft vs. Idiomatic Reconstruction}
\textbf{\color{TargetHighlight}[Defective Draft]}:\par\textit{For example, according to a recent survey, 92\% of factories that use robots become very rich and successful.}\par\vspace{4pt}
\textbf{\color{AcadNavy}[Idiomatic Reconstruction]}:\par\textit{For instance, when automotive assembly plants deploy collaborative robots for repetitive welding and heavy chassis lifting, human technicians can transition toward quality inspection and fault diagnostics, reducing workplace musculoskeletal injuries while raising line throughput.}\par\vspace{4pt}
\textbf{Surgical Analysis}: 不编造假百分比，而是通过“汽车总装厂部署协作机器人承担重复焊接与底盘搬运 $\to$ 技工转向质检与故障诊断 $\to$ 减少工伤并提升产线吞吐量”展现母语级具体说服力。
\end{ExampleBox}
\begin{TipBox}{Academic Trap \& Reviewer Perspective}
\textbf{Common Pitfall}: 举例要么假大空（For example, Jack studied hard and succeeded），要么编造一眼假的统计数字。
\end{TipBox}
\vspace{6pt}
\subsection{L022: Eliminating Tautology and Circular Explanations / 杜绝循环论证与同义反复（Eliminating Circular Reasoning / Tautology）}\label{sec:logic:L022}
\begin{ConceptBox}{Core Logic Principle \& Structural Formula}
\textbf{Category}: \texttt{Logical Fallacy Prevention (逻辑谬误防范)}\par\vspace{3pt}
\textbf{Core Principle}: 原因分句与结果分句不能只是换个近义词把同一件事说两遍（如 *The model generalizes well because it has strong generalization ability*）。because / since / as 之后必须提供处于更底层的**结构、物理或统计机制**。\par\vspace{4pt}
\textbf{Structural Formula}:\par\noindent\texttt{Claim: [High-Level Property A]. $\to$ Reason: Because [Low-Level Structural/Statistical Mechanism B], NOT [Synonym of Property A].}
\end{ConceptBox}
\begin{ExampleBox}{Defective Draft vs. Idiomatic Reconstruction}
\textbf{\color{TargetHighlight}[Defective Draft]}:\par\textit{Our policy is more adaptable to new objects because it possesses superior cross-object adaptability.}\par\vspace{4pt}
\textbf{\color{AcadNavy}[Idiomatic Reconstruction]}:\par\textit{Our policy adapts reliably to unseen object geometries because it conditions grasp proposals on local surface normals rather than global category templates.}\par\vspace{4pt}
\textbf{Surgical Analysis}: 将“具有跨物体适应性”替换为底层几何机制“以局部表面法向量而非全局类别模板作为抓取条件”，真正回答了 Why。
\end{ExampleBox}
\begin{TipBox}{Academic Trap \& Reviewer Perspective}
\textbf{Common Pitfall}: 经常写出“因为模型更鲁棒，所以它在噪声下表现更好”这种零信息增量的同义反复句。
\end{TipBox}
\vspace{6pt}
\subsection{L023: Moving Beyond False Dichotomies to Pareto Trade-offs / 破除二元对立谬误：展现多维权衡（Transcending False Dichotomies）}\label{sec:logic:L023}
\begin{ConceptBox}{Core Logic Principle \& Structural Formula}
\textbf{Category}: \texttt{Logical Fallacy Prevention (逻辑谬误防范)}\par\vspace{3pt}
\textbf{Core Principle}: 无论是讨论模型设计（精度 vs. 速度、生成多样性 vs. 保真度）还是社会议题（发展 AI vs. 保护就业、基础理论 vs. 工程应用），避免非黑即白的极端断言；用兼容性或帕累托改进框架（reconcile / bridge / balance X and Y by...）构建高阶论点。\par\vspace{4pt}
\textbf{Structural Formula}:\par\noindent\texttt{Rather than treating [Goal A] and [Goal B] as mutually exclusive, [Approach/Policy] reconciles both by [Mechanism C].}
\end{ConceptBox}
\begin{ExampleBox}{Defective Draft vs. Idiomatic Reconstruction}
\textbf{\color{TargetHighlight}[Defective Draft]}:\par\textit{Universities should stop teaching theoretical courses and only teach practical coding skills.}\par\vspace{4pt}
\textbf{\color{AcadNavy}[Idiomatic Reconstruction]}:\par\textit{Rather than treating mathematical foundations and engineering practice as mutually exclusive, computer science curricula can integrate both by embedding convex optimization and probability theory directly into hands-on machine learning studio projects.}\par\vspace{4pt}
\textbf{Surgical Analysis}: 用 Rather than treating A and B as mutually exclusive, ... reconciles/integrates both by... 展现成熟严密的学术思辨。
\end{ExampleBox}
\begin{TipBox}{Academic Trap \& Reviewer Perspective}
\textbf{Common Pitfall}: 非此即彼，为了夸大一端而把另一端贬得一无是处。
\end{TipBox}
\vspace{6pt}
\subsection{L024: Non-Stationary Drift and Multi-Scale Decomposition Logic / 时间序列非平稳性与多尺度解耦论证规范（Time Series Non-Stationarity Logic）}\label{sec:logic:L024}
\begin{ConceptBox}{Core Logic Principle \& Structural Formula}
\textbf{Category}: \texttt{Domain Discourse (时间序列专属论证)}\par\vspace{3pt}
\textbf{Core Principle}: 论述时间序列问题时，必须清晰区分两类矛盾：(1)  统计分布漂移（Distribution Shift：均值/方差随时间突变）与局部时序语义（Intrinsic Temporal Dependency）之间的矛盾；(2)  短期高频波动（High-frequency Fluctuations）与长期宏观趋势（Long-term Trend）的混叠干扰。\par\vspace{4pt}
\textbf{Structural Formula}:\par\noindent\texttt{Naive normalization removes [Non-stationary shift] at the expense of [Intrinsic scale information]. $\to$ Our reversible/decoupled module disentangles [Component A] from [Component B], restoring [Target Property].}
\end{ConceptBox}
\begin{ExampleBox}{Defective Draft vs. Idiomatic Reconstruction}
\textbf{\color{TargetHighlight}[Defective Draft]}:\par\textit{Real-world time series are very complicated and change all the time, so normal normalization makes predictions bad.}\par\vspace{4pt}
\textbf{\color{AcadNavy}[Idiomatic Reconstruction]}:\par\textit{While instance normalization mitigates non-stationary mean and variance shifts across lookback windows, stripping these statistics discards amplitude cues essential for forecasting peak electricity loads. Our adaptive de-stationary projection restores horizon-specific scale statistics during decoding, preserving both stationarity in attention computation and physical amplitude fidelity in the final output.}\par\vspace{4pt}
\textbf{Surgical Analysis}: 精准切中时间序列领域最核心的“去平稳化归一化抹除物理幅值信息 vs. 注意力层需要平稳分布”的张力。
\end{ExampleBox}
\begin{TipBox}{Academic Trap \& Reviewer Perspective}
\textbf{Common Pitfall}: 笼统说 time series data is very complex and hard to predict，说不清究竟是协变量偏移、频率混叠还是采样不规则。
\end{TipBox}
\vspace{6pt}
\subsection{L025: Embodied Perception-Action Loop and Sim-to-Real Logic / 具身智能感知-决策-控制闭环与 Sim-to-Real 论证规范（Embodied AI Closed-Loop Logic）}\label{sec:logic:L025}
\begin{ConceptBox}{Core Logic Principle \& Structural Formula}
\textbf{Category}: \texttt{Domain Discourse (具身智能专属论证)}\par\vspace{3pt}
\textbf{Core Principle}: 论述具身智能（Embodied AI）与机器人策略时，严禁将其等同于静态图像识别。必须突出**闭环反馈（Closed-loop feedback）**、**分布外状态累积偏移（Covariate shift in sequential decision-making）**、**接触动力学不确定性（Contact-rich dynamics）**以及**机载实时推理延迟约束（Onboard real-time latency）**。\par\vspace{4pt}
\textbf{Structural Formula}:\par\noindent\texttt{Unlike static perception where errors remain isolated, closed-loop manipulation compounds small action deviations into out-of-distribution states. $\to$ To maintain stability under contact uncertainty, our policy [Mechanism].}
\end{ConceptBox}
\begin{ExampleBox}{Defective Draft vs. Idiomatic Reconstruction}
\textbf{\color{TargetHighlight}[Defective Draft]}:\par\textit{The robot sees the image from the camera and predicts the action to pick up the cup accurately.}\par\vspace{4pt}
\textbf{\color{AcadNavy}[Idiomatic Reconstruction]}:\par\textit{Unlike static visual recognition where prediction errors remain isolated across samples, closed-loop robotic manipulation compounds minor trajectory deviations into unfamiliar contact states. To prevent this covariate drift during precision assembly, our visuomotor policy conditions short-horizon action chunks on real-time wrist force-torque feedback.}\par\vspace{4pt}
\textbf{Surgical Analysis}: 鲜明对比静态视觉与闭环操作中的误差累积（covariate drift），并引出视觉-力觉多模态短视野动作块控制机制。
\end{ExampleBox}
\begin{TipBox}{Academic Trap \& Reviewer Perspective}
\textbf{Common Pitfall}: 把机器人控制写得像普通分类任务（The robot looks at the image and outputs the right class of movement）。
\end{TipBox}
\vspace{6pt}
\subsection{L026: Fidelity, Mode Coverage, and Sampling Efficiency Trilemma / 生成模型三元权衡论证规范：生成质量、模式覆盖度与采样效率（Generative Trilemma Logic）}\label{sec:logic:L026}
\begin{ConceptBox}{Core Logic Principle \& Structural Formula}
\textbf{Category}: \texttt{Domain Discourse (生成模型专属论证)}\par\vspace{3pt}
\textbf{Core Principle}: 分析生成模型（Diffusion, Flow Matching, Autoregressive, GAN, VAE）时，段落逻辑应围绕生成三难困境（Trilemma）展开：(1)  样本保真度（Sample Fidelity / Temporal Realism）、(2)  分布模式覆盖度/多样性（Mode Coverage / Diversity，防止 Mode Collapse）、(3)  推理采样步数与速度（Inference NFE / Sampling Latency）。\par\vspace{4pt}
\textbf{Structural Formula}:\par\noindent\texttt{While [Model Family A] excels at [Fidelity/Diversity], it incurs [High NFE / Mode Collapse]. $\to$ By reformulating the generation trajectory via [Rectified Flow / Consistency Distillation], our approach preserves [Property 1] while reducing sampling steps from [N] to [k].}
\end{ConceptBox}
\begin{ExampleBox}{Defective Draft vs. Idiomatic Reconstruction}
\textbf{\color{TargetHighlight}[Defective Draft]}:\par\textit{Diffusion models generate very realistic data, but they are too slow. Our model makes them fast and still good.}\par\vspace{4pt}
\textbf{\color{AcadNavy}[Idiomatic Reconstruction]}:\par\textit{While denoising diffusion models capture multi-modal trajectory distributions without mode collapse, their reliance on tens of iterative solver steps precludes real-time control. By straightening the probability flow ODE via reflow distillation, our generator preserves multi-modal path diversity while compressing inference to a single forward pass.}\par\vspace{4pt}
\textbf{Surgical Analysis}: 精准使用 multi-modal trajectory distributions、without mode collapse、iterative solver steps、straightening the probability flow ODE 等领域高阶表达。
\end{ExampleBox}
\begin{TipBox}{Academic Trap \& Reviewer Perspective}
\textbf{Common Pitfall}: 只说 the generated samples look very good and realistic，缺乏分布覆盖度与采样效率维度的严谨坐标。
\end{TipBox}
\vspace{6pt}
\subsection{L027: Inter-Paragraph Handshake and Bridge Clauses / 段尾留钩与段首咬合：跨段落无缝过渡法则（Paragraph-to-Paragraph Handshake）}\label{sec:logic:L027}
\begin{ConceptBox}{Core Logic Principle \& Structural Formula}
\textbf{Category}: \texttt{Paragraph Transitions (段间衔接)}\par\vspace{3pt}
\textbf{Core Principle}: 整篇小论文或大作文不是几个孤立段落的拼盘。上一段的末句应为下一段埋下伏笔，或者下一段的首句用前置状语短语（Having established [Paragraph N's core], we now turn to [Paragraph N+1's core] / While [Paragraph N's factor] accounts for X, [Paragraph N+1's factor] governs Y）完成段间握手。\par\vspace{4pt}
\textbf{Structural Formula}:\par\noindent\texttt{End of Para A: ...resolving [Issue X] in isolation is insufficient without addressing [Issue Y]. $\to$ Start of Para B: To address this [Issue Y], we introduce [Module B]...}
\end{ConceptBox}
\begin{ExampleBox}{Defective Draft vs. Idiomatic Reconstruction}
\textbf{\color{TargetHighlight}[Defective Draft]}:\par\textit{[End of Para 1] ...thus extracting periodic features. [Start of Para 2] Next, we use a contrastive loss function to train the network.}\par\vspace{4pt}
\textbf{\color{AcadNavy}[Idiomatic Reconstruction]}:\par\textit{[End of Para 1] ...thus extracting multi-period representations from raw sensor streams. [Start of Para 2] Having isolated these periodic representations, we next align them across heterogeneous operating regimes via a regime-aware contrastive objective.}\par\vspace{4pt}
\textbf{Surgical Analysis}: 用 Having isolated these periodic representations, we next... 将上一段的输出直接作为下一段损失函数设计的输入，全篇浑然一体。
\end{ExampleBox}
\begin{TipBox}{Academic Trap \& Reviewer Perspective}
\textbf{Common Pitfall}: 段落之间毫无过渡，上一段讲完频域特征，下一段突然毫无征兆地讲损失函数，像看幻灯片跳页。
\end{TipBox}
\vspace{6pt}
\subsection{L028: Diagnosing and Bridging Missing Logical Premises / 诊断与修复逻辑跳跃（Bridging Logical Leaps / Missing Premises）}\label{sec:logic:L028}
\begin{ConceptBox}{Core Logic Principle \& Structural Formula}
\textbf{Category}: \texttt{Logical Fallacy Prevention (逻辑谬误防范)}\par\vspace{3pt}
\textbf{Core Principle}: 作者脑子里清楚 A $\to$ B $\to$ C，但写出来只写了 A，紧接着就 wrote Therefore, C。修改时必须站在审稿人/读者视角，把隐含的中间桥接前提 B（Missing Premise）显式补写出来。\par\vspace{4pt}
\textbf{Structural Formula}:\par\noindent\texttt{Premise A: [Observation]. $\to$ Missing Bridge B: [Because A implies Property B under Condition K], $\to$ Conclusion C: [Therefore, Design/Policy C is necessary].}
\end{ConceptBox}
\begin{ExampleBox}{Defective Draft vs. Idiomatic Reconstruction}
\textbf{\color{TargetHighlight}[Defective Draft]}:\par\textit{Industrial time series contain missing values at irregular intervals. Therefore, we use a continuous-time neural ODE encoder.}\par\vspace{4pt}
\textbf{\color{AcadNavy}[Idiomatic Reconstruction]}:\par\textit{Industrial sensor streams frequently suffer from irregular missing timestamps due to packet loss. Because standard discrete-step RNNs assume uniform sampling intervals and distort time-delta dynamics when inputs are missing, we model latent state evolution between arbitrary observation times using a continuous-time neural ODE.}\par\vspace{4pt}
\textbf{Surgical Analysis}: 补全了关键中间前提：“离散步长 RNN 默认等间隔采样，在缺失值下会扭曲时间间隔动力学”，从而让引入连续时间神经 ODE 变得顺理成章。
\end{ExampleBox}
\begin{TipBox}{Academic Trap \& Reviewer Perspective}
\textbf{Common Pitfall}: 滥用 Therefore / Thus / Hence 强行连接两句根本没有直接推出关系的话。
\end{TipBox}
\vspace{6pt}
\subsection{L029: Stripping LLM Rhetorical Padding at the Paragraph Level / 篇章级去 AI 味手术：砍掉空洞修辞外壳，直奔命题内核（Zero AI-ese Discourse Surgery）}\label{sec:logic:L029}
\begin{ConceptBox}{Core Logic Principle \& Structural Formula}
\textbf{Category}: \texttt{Style \& Anti-AI-ese (去 AI 化文风)}\par\vspace{3pt}
\textbf{Core Principle}: 大模型生成的段落有三大典型臭味：(1)  句首喜欢戴高帽（In the rapidly evolving landscape of... / It serves as a testament to...）；(2)  段中喜欢滥用排比与抽象隐喻（not only... but also..., tapestry, beacon, realm, paradigm shift, delve into, underscore）；(3)  段尾喜欢升华口号（paving the way for a brighter future）。修改时必须一刀砍掉所有修辞外壳，首句第一词直接落到具体研究对象或议题实体上。\par\vspace{4pt}
\textbf{Structural Formula}:\par\noindent\texttt{Delete: 'In the realm of X, it is undeniable that Y plays a pivotal role in spearheading Z.' $\to$ Write: '[Concrete Entity Y] governs [Specific Mechanism Z] by [Measurable Action].'}
\end{ConceptBox}
\begin{ExampleBox}{Defective Draft vs. Idiomatic Reconstruction}
\textbf{\color{TargetHighlight}[Defective Draft]}:\par\textit{In the rapidly evolving landscape of artificial intelligence, time series forecasting stands as a cornerstone that plays a pivotal role in spearheading industrial transformation, serving as a testament to data-driven decision-making.}\par\vspace{4pt}
\textbf{\color{AcadNavy}[Idiomatic Reconstruction]}:\par\textit{Accurate multi-horizon time series forecasting enables power grid operators to balance renewable generation against fluctuating consumer demand.}\par\vspace{4pt}
\textbf{Surgical Analysis}: 砍掉 32 个词的 AI 废话外壳，用 16 个词的紧凑主谓宾直接锁定“多步时间序列预测帮电网调度员平衡新能源发电与负荷波动”。
\end{ExampleBox}
\begin{TipBox}{Academic Trap \& Reviewer Perspective}
\textbf{Common Pitfall}: 误以为辞藻越华丽、排比越工整、抽象大词越多，英语就越高级。
\end{TipBox}
\vspace{6pt}
\subsection{L030: Seven-Point Post-Drafting Paragraph Audit / 成段后七维自检清单（The 7-Point Paragraph Audit Checklist）}\label{sec:logic:L030}
\begin{ConceptBox}{Core Logic Principle \& Structural Formula}
\textbf{Category}: \texttt{Quality Assurance (段落自检闭环)}\par\vspace{3pt}
\textbf{Core Principle}: 写完或改完任何一个英语段落，必须依次用七把尺子过筛：(1)  首句是否一眼看清本段唯一核心主张？(2)  每句主语是否控制在 6--8 词以内（无头重脚轻）？(3)  句与句之间是否遵循“旧信息在前、新信息在后”？(4)  所有的 This/Such 后是否都跟了精准汇总名词？(5)  是否彻底清除了口语词与 AI 套话？(6)  所有的程度主张是否都有具体指标、基线或微观案例支撑？(7)  主动与被动语态的切换是否服务于信息焦点？\par\vspace{4pt}
\textbf{Structural Formula}:\par\noindent\texttt{Audit(Para) = Unity(1) + ShortSubject(2) + GivenToNew(3) + SummaryNoun(4) + ZeroAI\_ZeroColloquial(5) + ConcreteMetricOrCase(6) + PurposefulVoice(7).}
\end{ConceptBox}
\begin{ExampleBox}{Defective Draft vs. Idiomatic Reconstruction}
\textbf{\color{TargetHighlight}[Defective Draft]}:\par\textit{Checking only spelling and grammar while leaving a 22-word subject, bare 'This', and unsubstantiated 'very good performance' untouched.}\par\vspace{4pt}
\textbf{\color{AcadNavy}[Idiomatic Reconstruction]}:\par\textit{Applying the 7-point audit systematically: compressing the 22-word subject to 4 words, replacing bare 'This' with 'This phase lag', and replacing 'very good' with 'a 14.6\% MSE reduction on Weather'.}\par\vspace{4pt}
\textbf{Surgical Analysis}: 将抽象的“改作文/改论文”转化为 7 个可逐项打勾的工程化核验动作。
\end{ExampleBox}
\begin{TipBox}{Academic Trap \& Reviewer Perspective}
\textbf{Common Pitfall}: 写完只检查拼写和单复数，不检查段落骨架、信息流与论证闭环。
\end{TipBox}
\vspace{6pt}